\documentclass[11pt, a4paper]{ctexart}

\usepackage{assignment_style}

\begin{document}

% =========================================
% 1
% =========================================
\begin{problem}{1. 用幂法计算下列矩阵的按模最大特征值及相应特征向量，取 $x^{(0)}=(1,0,0)^{\mathrm{T}}$。}
    \[
        \text{(1)}\quad
        A_1 =
        \begin{pmatrix}
            -4 & 14 & 0 \\
            -5 & 13 & 0 \\
            -1 & 0 & 2
        \end{pmatrix},
        \qquad
        \text{(2)}\quad
        A_2 =
        \begin{pmatrix}
            4 & -1 & 1 \\
            -1 & 3 & -2 \\
            1 & -2 & 3
        \end{pmatrix}.
    \]
\end{problem}

\begin{solution}
    采用归一化幂法。设当前向量为 $x$，取
    \[
        r = \arg\max_i |x_i|,
        \qquad
        x \leftarrow \frac{x}{x_r},
    \]
    然后计算
    \[
        x \leftarrow Ax,\qquad
        \lambda_{\text{new}} = x_r,\quad r = \arg\max_i |x_i|.
    \]
    输出近似特征向量 $x/\lambda_{\text{new}}$ 与近似特征值 $\lambda_{\text{new}}$；若
    \[
        |\lambda_{\text{new}}-\lambda_{\text{old}}|<\varepsilon,
    \]
    则停止迭代。

    \textbf{(1)} 对
    \[
        A_1 =
        \begin{pmatrix}
            -4 & 14 & 0 \\
            -5 & 13 & 0 \\
            -1 & 0 & 2
        \end{pmatrix},
    \]
    迭代结果为
    \[
        \renewcommand{\arraystretch}{1.15}
        \begin{array}{c|rrrr}
            k & x_1 & x_2 & x_3 & \lambda \\ \hline
            1 & 1 & 0 & 0 & 0 \\
            2 & 0.8 & 1 & 0.2 & -5 \\
            3 & 1 & 0.8333 & -0.0370 & 10.800 \\
            \cdots & \cdots & \cdots & \cdots & \cdots \\
            10 & 1 & 0.7146 & -0.2490 & 6.0084
        \end{array}
    \]
    因而按模最大特征值近似为
    \[
        \lambda \approx 6.0084,
        \qquad
        x \approx (1,\,0.7146,\,-0.2490)^{\mathrm{T}}.
    \]

    \textbf{(2)} 对
    \[
        A_2 =
        \begin{pmatrix}
            4 & -1 & 1 \\
            -1 & 3 & -2 \\
            1 & -2 & 3
        \end{pmatrix},
    \]
    迭代结果为
    \[
        \renewcommand{\arraystretch}{1.15}
        \begin{array}{c|rrrr}
            k & x_1 & x_2 & x_3 & \lambda \\ \hline
            1 & 1 & 0 & 0 & 0 \\
            2 & 1 & -0.25 & 0.25 & 4 \\
            3 & 1 & -0.5 & 0.5 & 4.5 \\
            \cdots & \cdots & \cdots & \cdots & \cdots \\
            10 & 1 & -0.9971 & 0.9971 & 5.9883
        \end{array}
    \]
    因而按模最大特征值近似为
    \[
        \lambda \approx 5.9883,
        \qquad
        x \approx (1,\,-0.9971,\,0.9971)^{\mathrm{T}}.
    \]
\end{solution}


% =========================================
% 3
% =========================================
\begin{problem}{3. 求矩阵最接近 $9.6$ 的特征值及对应特征向量。}
    \[
        A =
        \begin{pmatrix}
            1 & 2 & 3 \\
            2 & 3 & 4 \\
            3 & 4 & 5
        \end{pmatrix}.
    \]
\end{problem}

\begin{solution}
    取位移 $\lambda_0=9.6$，问题转化为求
    \[
        (A-\lambda_0 I)^{-1}
    \]
    的按模最大特征值及对应特征向量。

    反幂法迭代步骤为：对 $A-\lambda_0 I$ 作 $LU$ 分解，解
    \[
        (A-\lambda_0 I)y = x,
    \]
    并用 $y$ 的绝对值最大分量归一化。设反幂法得到的特征值近似为 $\lambda^\ast$，则原矩阵 $A$ 的对应特征值为
    \[
        \lambda = \frac{1}{\lambda^\ast}+\lambda_0.
    \]

    迭代结果为
    \[
        \renewcommand{\arraystretch}{1.15}
        \begin{array}{c|rrrr}
            k & x_1 & x_2 & x_3 & \dfrac{1}{\lambda^\ast}+\lambda_0 \\ \hline
            1 & 1 & 0 & 0 & 19.2 \\
            2 & 0.5165 & 0.7626 & 1 & 9.6829 \\
            3 & 0.5247 & 0.7623 & 1 & 9.6235 \\
            \cdots & \cdots & \cdots & \cdots & \cdots \\
            8 & 0.5247 & 0.7623 & 1 & 9.6235
        \end{array}
    \]
    因而最接近 $9.6$ 的特征值近似为
    \[
        \lambda \approx 9.6235,
        \qquad
        x \approx (0.5247,\,0.7623,\,1)^{\mathrm{T}}.
    \]
\end{solution}


% =========================================
% 4
% =========================================
\begin{problem}{4. 用反幂法计算矩阵}
    \[
        A =
        \begin{pmatrix}
            2 & 1 & 0 \\
            1 & 3 & 1 \\
            0 & 1 & 4
        \end{pmatrix}
    \]
    相应于特征值 $1.2679$ 的特征向量。
\end{problem}

\begin{solution}
    令 $\lambda_0=1.2679$，仍对 $(A-\lambda_0 I)^{-1}$ 使用反幂法。
    迭代结果为
    \[
        \renewcommand{\arraystretch}{1.15}
        \begin{array}{c|rrrr}
            k & x_1 & x_2 & x_3 & \dfrac{1}{\lambda^\ast}+\lambda_0 \\ \hline
            1 & 1 & 0 & 0 & 2.5358 \\
            2 & 1 & -0.7320 & 0.2679 & 1.2680 \\
            3 & 1 & -0.7321 & 0.2679 & 1.2679 \\
            \cdots & \cdots & \cdots & \cdots & \cdots \\
            5 & 1 & -0.7321 & 0.2679 & 1.2679
        \end{array}
    \]
    故相应于特征值 $1.2679$ 的特征向量可取
    \[
        x \approx (1,\,-0.7321,\,0.2679)^{\mathrm{T}}.
    \]
\end{solution}

\end{document}
